\documentclass[10pt]{article}
\usepackage[letterpaper,margin=0.62in]{geometry}
\usepackage{amsmath,amssymb,mathtools}
\usepackage{array,booktabs,multicol}
\usepackage{enumitem}
\usepackage{xcolor}
\usepackage{hyperref}
\usepackage{listings}
\usepackage{titlesec}
\usepackage{microtype}

\hypersetup{hidelinks}
\setlength{\parindent}{0pt}
\setlength{\parskip}{4pt}
\setlist[itemize]{leftmargin=*,topsep=2pt,itemsep=1pt}
\setlist[enumerate]{leftmargin=*,topsep=2pt,itemsep=1pt}

\titleformat{\section}{\Large\bfseries}{\thesection}{0.5em}{}
\titleformat{\subsection}{\large\bfseries}{\thesubsection}{0.5em}{}
\titleformat{\subsubsection}{\normalsize\bfseries}{\thesubsubsection}{0.5em}{}

\lstset{
  basicstyle=\ttfamily\footnotesize,
  columns=fullflexible,
  keepspaces=true,
  frame=single,
  breaklines=true,
  showstringspaces=false
}

\newcommand{\bigO}{O}
\newcommand{\SA}{\operatorname{SA}}
\newcommand{\LCP}{\operatorname{LCP}}
\newcommand{\rankv}{\operatorname{rank}}
\newcommand{\lca}{\operatorname{lca}}
\newcommand{\dist}{\operatorname{dist}}
\newcommand{\RMQ}{\operatorname{RMQ}}
\newcommand{\pred}{\operatorname{pred}}
\newcommand{\succf}{\operatorname{succ}}

\newcommand{\boxhead}[1]{\vspace{3pt}\textbf{#1.}}

\title{\vspace{-1.0cm}Ejercicios Flash para el Final de EDA\\
\large Guia formal de teoria aplicada, soluciones y complejidades}
\author{}
\date{}

\begin{document}
\maketitle
\vspace{-0.7cm}

\section*{Como usar esta guia}
Este documento reemplaza la version anterior de \texttt{ejercicios\_flash.pdf}. Esta escrito para responder
problemas teoricos tipo Racso: no basta nombrar una estructura; hay que justificar por que aplica, que guarda,
como responde, por que es correcto y cuanto cuesta.

\boxhead{Orden de estudio si falta tiempo}
\begin{enumerate}
\item Lee primero la Seccion 1: plantilla de respuesta y mapa de decision.
\item Luego asegura las tres preguntas de la lista del profesor: Secciones 2, 3 y 4.
\item Despues strings: suffix array, LCP y aplicaciones (Secciones 6--9).
\item Finalmente dynamic e integer: rollback DSU, Euler-tour/link-cut, vEB/Y-fast/radix (Secciones 12--15).
\end{enumerate}

\boxhead{Formato que debe sonar en el examen}
``El problema es estatico/dinamico. La estructura correcta es X porque transforma la consulta en Y.
Preproceso guardando Z. Para cada query hago estos pasos. La correctitud se basa en este invariante/propiedad.
La complejidad temporal es ... y la memoria es ...''.

\tableofcontents
\newpage

\section{Plantilla formal y mapa de decision}

\subsection{Plantilla de respuesta para cualquier problema}
Para cada ejercicio, intenta llenar esta plantilla en papel:
\begin{enumerate}
\item \textbf{Objeto}: arreglo, string, arbol, bosque dinamico, grafo general, conjunto de enteros.
\item \textbf{Tipo}: estatico, incremental, decremental, fully dynamic, online u offline.
\item \textbf{Estructura}: sparse table, suffix array, LCA, HLD, centroid decomposition, rollback DSU,
Euler-tour tree, link-cut tree, vEB/Y-fast, radix sort.
\item \textbf{Que guarda}: por ejemplo, \(\LCP\), ancestors \(2^j\), maximo en salto, componentes DSU,
intervalos de vida de aristas, prefijos de un trie binario.
\item \textbf{Query/update}: traducir la pregunta al lenguaje de la estructura.
\item \textbf{Correctitud}: explicar el invariante que hace que la traduccion sea exacta.
\item \textbf{Complejidad}: preprocesamiento, query/update y memoria.
\end{enumerate}

\subsection{Mapa de decision}
\begin{center}
\begin{tabular}{p{0.38\linewidth}p{0.26\linewidth}p{0.25\linewidth}}
\toprule
Frase del enunciado & Estructura & Complejidad tipica \\
\midrule
Arreglo fijo, min/max/gcd en rangos & Sparse table & build \(\bigO(n\log n)\), query \(\bigO(1)\) \\
Camino en arbol fijo & LCA / HLD & LCA \(\bigO(\log n)\), HLD \(\bigO(\log^2 n)\) \\
Subarbol en arbol fijo & Euler tour timestamps & query/update como intervalo \\
Muchos patrones en texto fijo & Suffix array o Aho-Corasick & depende de patrones \\
Comparar/ordenar subcadenas & SA + LCP + RMQ & comparacion \(\bigO(1)\) \\
Contar caminos en arbol con restriccion de distancia & Centroid decomposition & \(\bigO(n\log^2 n)\) \\
Add/delete edges offline & Segment tree tiempo + rollback DSU & \(\bigO(m\log q)\) uniones \\
Bosque dinamico con path aggregate & Link-cut tree & \(\bigO(\log n)\) amortizado \\
Bosque dinamico con componente/subtree aggregate & Euler-tour tree & \(\bigO(\log n)\) \\
Predecessor de enteros & vEB / Y-fast / Fusion & \(\bigO(\log w)\), \(\bigO(\log_w n)\) \\
Ordenar enteros & Counting/radix/signature & puede romper \(n\log n\) \\
\bottomrule
\end{tabular}
\end{center}

\newpage

\section{Lista profesor 1: subsecuencia en un camino de arbol}

\subsection{Enunciado tipo}
Se tiene un arbol con \(N\) vertices. Cada vertice contiene una letra minuscula. En una consulta se dan
\(U,V\) y una cadena \(S\). Sea \(T\) la cadena obtenida al escribir, en orden, las letras del camino simple
de \(U\) a \(V\). Se debe decidir si \(S\) es subsecuencia de \(T\).

\subsection{Reconocimiento}
El arbol es fijo y las queries son sobre caminos. La palabra clave es \textbf{camino en arbol estatico}.
La estructura natural es HLD, porque descompone cualquier camino en \(\bigO(\log N)\) segmentos de cadenas
pesadas. Como el alfabeto es pequeno, se puede guardar para cada letra las posiciones donde aparece en el
arreglo lineal de HLD.

\subsection{Estructuras que se guardan}
\begin{itemize}
\item HLD: \(parent[v]\), \(depth[v]\), \(head[v]\), \(pos[v]\), \(out[v]\).
\item Arreglo base \(B[pos[v]] =\) letra del vertice \(v\).
\item Para cada letra \(c\in\{a,\ldots,z\}\), lista ordenada \(L_c=\{pos[v] : B[pos[v]]=c\}\).
\end{itemize}

\subsection{Algoritmo}
Primero se descompone el camino \(U\to V\) en segmentos sobre el arreglo base, respetando el orden real
del camino. Hay dos partes:
\[
U \to \lca(U,V) \quad\text{y}\quad \lca(U,V)\to V.
\]
La primera parte se recorre hacia arriba; dentro de una cadena HLD eso puede significar posiciones
decrecientes. La segunda parte se recorre hacia abajo; se guardan sus segmentos y se invierten al final.

Luego se consume la cadena \(S\) de izquierda a derecha. Si el caracter actual es \(S[i]\), se busca en
la lista \(L_{S[i]}\) la primera aparicion valida dentro del segmento actual y despues de la posicion ya usada.
En segmentos crecientes se usa \texttt{lower\_bound}; en segmentos decrecientes se usa \texttt{upper\_bound}
y se retrocede.

\begin{lstlisting}
query(U,V,S):
  segments = ordered_HLD_segments(U,V)  // each segment has [l,r] and direction
  i = 0
  for seg in segments:
    cur = seg.start
    while i < |S|:
      p = first_position_of_char_inside_segment_after(cur, S[i], seg)
      if p does not exist: break
      cur = next_position_after(p, seg.direction)
      i++
  return i == |S|
\end{lstlisting}

\subsection{Correctitud}
HLD no altera el camino; solo lo parte en segmentos contiguos de cadenas pesadas. Por tanto, concatenar
los segmentos en el orden correcto produce exactamente \(T\). Dentro de un segmento, las listas \(L_c\)
permiten encontrar la primera aparicion de un caracter que todavia puede usarse.

El greedy de tomar la primera aparicion valida es correcto para subsecuencias: si existe una solucion que usa
una aparicion posterior del caracter actual, reemplazarla por una aparicion anterior no invalida los caracteres
restantes, porque deja un sufijo del camino igual o mas grande para continuar. Por induccion sobre los caracteres
de \(S\), si el algoritmo falla en algun caracter, ninguna subsecuencia completa existe.

\subsection{Complejidad}
Preprocesamiento HLD: \(\bigO(N)\). Construir listas por caracter: \(\bigO(N)\) tiempo y espacio.
Una query usa \(\bigO(\log N)\) segmentos. Cada caracter consumido hace una busqueda binaria en una lista,
costando \(\bigO(\log N)\). Por tanto:
\[
T_{\text{query}}=\bigO((|S|+\log N)\log N), \qquad
M=\bigO(N).
\]
Si se optimizan las listas por cadena pesada, puede bajarse constante, pero esta respuesta ya es claramente
menor que recorrer todo el camino para cada consulta.

\section{Lista profesor 2: ordenar nodos de un trie}

\subsection{Enunciado tipo}
Dado un trie \(T\), cada nodo \(u\) representa la cadena formada por las etiquetas de las aristas desde la raiz
hasta \(u\). Se pide ordenar los nodos por la comparacion lexicografica de sus cadenas en complejidad menor
que \(\bigO(n^2)\).

\subsection{Idea}
En orden lexicografico, si una palabra \(A\) es prefijo propio de \(B\), entonces \(A<B\). Por eso un nodo debe
salir antes que todos sus descendientes. Entre dos hijos de un mismo nodo, el caracter menor debe procesarse
primero. Esto sugiere DFS preorder con hijos ordenados.

\subsection{Algoritmo}
\begin{lstlisting}
lexicographic_trie_order(root):
  answer = []
  dfs(root)

dfs(u):
  append u to answer
  for character c in increasing order:
    if child[u][c] exists:
      dfs(child[u][c])
\end{lstlisting}

Si la raiz representa la cadena vacia y el problema pide ordenar todos los nodos, la raiz va primero. Si solo se
piden cadenas no vacias, se omite la raiz del arreglo final.

\subsection{Correctitud}
Sea \(u\) un nodo y \(v\) un descendiente propio. La cadena de \(u\) es prefijo propio de la cadena de \(v\);
por definicion lexicografica, \(str(u)<str(v)\). Por eso visitar \(u\) antes que su subarbol es correcto.

Ahora considera dos nodos que pertenecen a subarboles de hijos distintos \(a\) y \(b\) de un mismo nodo, con
etiquetas \(c_a<c_b\). Sus cadenas comparten el prefijo del padre y difieren por primera vez en esas etiquetas,
por lo que toda cadena del subarbol de \(a\) es menor que toda cadena del subarbol de \(b\). Visitar hijos en
orden creciente respeta ese criterio. Por induccion sobre la profundidad del trie, el DFS produce el orden
lexicografico global.

\subsection{Complejidad}
Si el alfabeto es fijo, por ejemplo \(\sigma=26\), iterar todos los caracteres en cada nodo cuesta:
\[
\bigO(n\sigma)=\bigO(n).
\]
La memoria adicional es \(\bigO(h)\) por recursion, donde \(h\) es la altura del trie, o \(\bigO(n)\) en el peor caso.
Si los hijos estan guardados en listas desordenadas, se ordenan por caracter y el costo total es
\(\sum_u \deg(u)\log \deg(u)\le \bigO(n\log n)\), todavia menor que \(\bigO(n^2)\).

\section{Lista profesor 3: cantidad de caminos con peso \(\le C\)}

\subsection{Enunciado tipo}
Dado un arbol ponderado \(T\), calcular la cantidad de caminos \(u\rightsquigarrow v\) cuyo peso total sea
menor o igual que una constante \(C\). El algoritmo debe ser menor que \(\bigO(n^2)\).

\subsection{Por que no basta LCA/HLD}
LCA calcula el peso de un camino especifico en \(\bigO(\log n)\). Pero aqui se piden muchos caminos: en el peor
caso hay \(\Theta(n^2)\) pares \((u,v)\). Necesitamos contar pares sin enumerarlos. La estructura clasica para
contar caminos con restricciones de distancia en arboles es \textbf{centroid decomposition}.

\subsection{Idea de centroid decomposition}
Un centroid \(c\) es un nodo cuya remocion deja componentes de tamano a lo mucho \(n/2\). En cualquier nivel
de la recursion:
\begin{itemize}
\item contamos todos los caminos que pasan por el centroid \(c\);
\item restamos los pares que caen completamente dentro del mismo subarbol hijo, porque esos no pasan por \(c\);
\item recursamos en cada componente restante.
\end{itemize}

\subsection{Como contar pares con suma \(\le C\)}
Recolecta todas las distancias desde \(c\) a los nodos de la componente:
\[
D=\{\dist(c,x): x \text{ en la componente}\}.
\]
Ordena \(D\). Con dos punteros:
\begin{lstlisting}
count_pairs(D,C):
  sort(D)
  l = 0; r = |D|-1; ans = 0
  while l < r:
    if D[l] + D[r] <= C:
      ans += r-l
      l++
    else:
      r--
  return ans
\end{lstlisting}
Esto cuenta pares no ordenados \(\{u,v\}\) con \(\dist(c,u)+\dist(c,v)\le C\).

\subsection{Algoritmo completo}
\begin{lstlisting}
solve(component):
  c = centroid(component)
  mark c as removed
  all = [0]  // distance from c to c
  for each neighbor subtree S of c not removed:
    temp = distances from c to nodes in S
    answer -= count_pairs(temp, C)
    append temp to all
  answer += count_pairs(all, C)
  for each remaining component S:
    solve(S)
\end{lstlisting}

\subsection{Correctitud}
Todo camino \(u\rightsquigarrow v\) tiene un primer centroid en la recursion que separa a \(u\) y \(v\) en
componentes distintas, o coincide con uno de los extremos. En ese nivel, el camino pasa por el centroid \(c\),
por lo que su peso es:
\[
\dist(u,v)=\dist(c,u)+\dist(c,v).
\]
El conteo global \(\texttt{count\_pairs(all,C)}\) cuenta ese par. Si ambos extremos estaban en el mismo subarbol
de \(c\), el camino no pasa por \(c\), por eso se resta en \(\texttt{count\_pairs(temp,C)}\) y se deja para la
recursion dentro de ese subarbol. Asi, cada camino valido se cuenta exactamente una vez.

\subsection{Complejidad}
Cada nivel de centroid decomposition reduce el tamano de la componente al menos a la mitad, por lo que hay
\(\bigO(\log n)\) niveles por nodo. Si en cada nivel se ordenan listas, una cota simple es:
\[
T(n)=\bigO(n\log^2 n), \qquad M(n)=\bigO(n).
\]
Esto cumple claramente menor que \(\bigO(n^2)\). Con tecnicas de merge o manteniendo listas ordenadas puede
mejorarse a \(\bigO(n\log n)\), pero para respuesta de examen la version \(\bigO(n\log^2 n)\) es defendible.

\newpage

\section{Metricas: bloque de emergencia de la lista}

\subsection{Definicion}
Una funcion \(d:X\times X\to \mathbb{R}\) es metrica si cumple:
\begin{enumerate}
\item No negatividad: \(d(x,y)\ge 0\).
\item Identidad: \(d(x,y)=0 \Leftrightarrow x=y\).
\item Simetria: \(d(x,y)=d(y,x)\).
\item Desigualdad triangular: \(d(x,z)\le d(x,y)+d(y,z)\).
\end{enumerate}

\subsection{Suma de metricas}
Si \(d_1,d_2\) son metricas y \(d_3(x,y)=d_1(x,y)+d_2(x,y)\), entonces \(d_3\) es metrica.
La no negatividad y simetria son inmediatas. Para identidad, si \(d_3(x,y)=0\), como ambos terminos son no
negativos, \(d_1(x,y)=d_2(x,y)=0\), luego \(x=y\). Para triangular:
\[
d_3(x,z)=d_1(x,z)+d_2(x,z)
\le d_1(x,y)+d_1(y,z)+d_2(x,y)+d_2(y,z)
=d_3(x,y)+d_3(y,z).
\]

\subsection{Como verificar funciones rapidamente}
\begin{itemize}
\item \((x-y)^2\) no es metrica: \(d(0,2)=4\), pero \(d(0,1)+d(1,2)=1+1=2\).
\item \(|x-y|+|f(x)-f(y)|\) suele ser metrica si incluye \(|x-y|\): el segundo termino es una pseudometrica y
el primero garantiza identidad.
\item \(L_1\), \(L_2\) y \(L_\infty\) son metricas en \(\mathbb{R}^d\).
\item Pesos positivos de \(L_1\), por ejemplo \(2|x_1-y_1|+3|x_2-y_2|\), son metricas.
\item \(\sqrt{|x-y|}\) es metrica porque \(\sqrt{a+b}\le \sqrt a+\sqrt b\).
\item \(\min\{1,|x-y|\}\) es metrica: truncar una metrica por \(1\) preserva triangular.
\item \(\ln(1+|x-y|)\) es metrica porque
\[
\ln(1+a+b)\le \ln((1+a)(1+b))=\ln(1+a)+\ln(1+b).
\]
\item \(\frac{|x-y|}{1+|x-y|}\) es metrica porque \(f(t)=t/(1+t)\) es creciente y subaditiva.
\end{itemize}

\newpage

\section{Static RMQ y Sparse Table}

\subsection{Contexto}
RMQ significa Range Minimum Query. Dado un arreglo fijo \(A\), una consulta pide el minimo en un rango.
Si no hay updates, una sparse table es ideal para operaciones idempotentes:
\[
\min(x,x)=x,\quad \max(x,x)=x,\quad \gcd(x,x)=x.
\]
No sirve igual para suma con dos bloques solapados, porque \(x+x\neq x\).

\subsection{Definicion}
\[
st[k][i]=\min(A[i],A[i+1],\ldots,A[i+2^k-1]).
\]
La transicion parte el bloque en dos mitades:
\[
st[k][i]=\min(st[k-1][i],\,st[k-1][i+2^{k-1}]).
\]
Para una consulta cerrada \([l,r]\):
\[
len=r-l+1,\qquad k=\lfloor \log_2 len\rfloor,
\]
\[
\RMQ(l,r)=\min(st[k][l],\,st[k][r-2^k+1]).
\]
Para una consulta medio abierta \([l,r)\), se usa \(len=r-l\) y:
\[
\RMQ(l,r)=\min(st[k][l],\,st[k][r-2^k]).
\]

\subsection{Pseudocodigo}
\begin{lstlisting}
build(A):
  lg[1] = 0
  for i = 2..n: lg[i] = lg[i/2] + 1
  for i = 0..n-1: st[0][i] = A[i]
  for k = 1..LG:
    for i = 0; i + 2^k <= n; i++:
      st[k][i] = min(st[k-1][i], st[k-1][i+2^(k-1)])

query(l,r): // closed interval
  len = r-l+1
  k = lg[len]
  return min(st[k][l], st[k][r-2^k+1])
\end{lstlisting}

\subsection{Correctitud}
Por induccion en \(k\), \(st[k][i]\) guarda el minimo del bloque de longitud \(2^k\). Para \(k=0\), el bloque
tiene un solo elemento. Para \(k>0\), el bloque se divide en dos bloques correctos de longitud \(2^{k-1}\).

En una query, los dos bloques de longitud \(2^k\) elegidos cubren todo el rango. Pueden solaparse, pero el minimo
no cambia por contar un elemento dos veces. Por tanto, la respuesta es correcta.

\subsection{Complejidad}
\[
T_{\text{build}}=\bigO(n\log n),\qquad T_{\text{query}}=\bigO(1),\qquad M=\bigO(n\log n).
\]

\newpage

\section{Suffix Array, Rank, LCP y Kasai}

\subsection{Definiciones basicas}
Un \textbf{sufijo} de \(S\) es \(S[i..n-1]\). Un \textbf{prefijo} es una parte inicial de una cadena.
El suffix array \(\SA\) guarda los indices de los sufijos en orden lexicografico:
\[
S[\SA[0]..] < S[\SA[1]..] < \cdots < S[\SA[n-1]..].
\]
El arreglo inverso \(\rankv\) cumple:
\[
\rankv[i]=p \quad \Longleftrightarrow \quad \SA[p]=i.
\]
El arreglo LCP guarda:
\[
\LCP[p]=\operatorname{lcp}(S[\SA[p-1]..],S[\SA[p]..]).
\]

\subsection{Ejemplo: banana}
Los sufijos de \texttt{banana} ordenados son:
\[
\texttt{a}(5),\ \texttt{ana}(3),\ \texttt{anana}(1),\ \texttt{banana}(0),\ \texttt{na}(4),\ \texttt{nana}(2).
\]
Entonces:
\[
\SA=[5,3,1,0,4,2].
\]
Como el sufijo que empieza en \(1\) es \texttt{anana} y ocupa la posicion \(2\), \(\rankv[1]=2\).

\subsection{Manber-Myers}
La idea es ordenar por prefijos de longitud \(1,2,4,8,\ldots\). Si ya conocemos el rank de prefijos de longitud
\(len\), entonces el prefijo de longitud \(2len\) del sufijo \(i\) se representa por:
\[
(\rankv[i],\rankv[i+len]).
\]
Ordenar estos pares ordena los prefijos de longitud \(2len\). Pares iguales reciben el mismo nuevo rank.

\subsection{Kasai}
Kasai construye \(\LCP\) en tiempo lineal despues de tener \(\SA\) y \(\rankv\). Para cada posicion \(i\), mira
el sufijo anterior en el orden lexicografico:
\[
j=\SA[\rankv[i]-1].
\]
Si el LCP anterior era \(h\), al pasar de \(i\) a \(i+1\) el LCP baja como mucho en \(1\), por lo que se empieza
desde \(\max(h-1,0)\) y se avanza mientras coincidan caracteres.

\subsection{Formula clave con RMQ}
Si \(a=\rankv[i]<b=\rankv[j]\), entonces:
\[
\operatorname{lcp}(i,j)=\min(\LCP[a+1],\LCP[a+2],\ldots,\LCP[b]).
\]
Esto convierte LCP entre dos sufijos arbitrarios en una consulta RMQ.

\subsection{Complejidad}
Manber-Myers con sort comun cuesta \(\bigO(n\log^2 n)\); optimizado con counting/radix por ranks cuesta
\(\bigO(n\log n)\). Kasai cuesta \(\bigO(n)\). El RMQ sobre LCP cuesta \(\bigO(n\log n)\) memoria y build,
y responde en \(\bigO(1)\).

\section{Aplicaciones obligatorias de Suffix Array}

\subsection{Buscar y contar ocurrencias de un patron}
Todos los sufijos que empiezan con un patron \(P\) forman un intervalo contiguo en \(\SA\). Se hacen dos busquedas
binarias:
\begin{itemize}
\item \(L\): primer sufijo lexicograficamente \(\ge P\).
\item \(R\): primer sufijo que ya no tiene prefijo \(P\).
\end{itemize}
Si \(L=R\), el patron no aparece. Si no:
\[
\#occ(P)=R-L,\qquad posiciones=\SA[L],\ldots,\SA[R-1].
\]
Complejidad simple por patron: \(\bigO(|P|\log n)\). Con comparaciones aceleradas por LCP puede mejorarse,
pero para examen la idea central es el intervalo.

\subsection{Subcadenas distintas}
Cada subcadena es prefijo de algun sufijo. El sufijo \(\SA[i]\) aporta \(n-\SA[i]\) prefijos, pero los primeros
\(\LCP[i]\) ya aparecieron en el sufijo anterior. Por tanto:
\[
\#\text{subcadenas distintas}
=\sum_i \big((n-\SA[i])-\LCP[i]\big)
=\frac{n(n+1)}2-\sum_i\LCP[i].
\]

\subsection{Longest Common Substring}
Para dos cadenas \(A,B\), construir:
\[
T=A\#B\$,
\]
donde \(\#,\$\) no aparecen en las cadenas. En el suffix array de \(T\), se revisan vecinos de origen distinto.
El mayor LCP entre vecinos de distinto origen es una subcadena comun maxima.

Correctitud: cualquier subcadena comun es prefijo de un sufijo de \(A\) y de un sufijo de \(B\). En el bloque de
sufijos que comparten ese prefijo debe existir algun par vecino de origen distinto; por eso revisar vecinos basta.

\subsection{Comparar y ordenar subcadenas}
Para comparar \(A=S[l_1..r_1]\) y \(B=S[l_2..r_2]\), sea:
\[
x=\min(\operatorname{lcp}(l_1,l_2), |A|, |B|).
\]
Si \(x<\min(|A|,|B|)\), decide el primer caracter distinto:
\[
S[l_1+x] \quad \text{vs.}\quad S[l_2+x].
\]
Si \(x=\min(|A|,|B|)\), una es prefijo de la otra y gana la mas corta. Si son identicas y el statement pide
desempate por \((l,r)\), se compara el par original. Con SA+LCP+RMQ, cada comparacion cuesta \(\bigO(1)\);
ordenar \(q\) subcadenas cuesta \(\bigO(q\log q)\).

\subsection{Bordes de todas las subcadenas}
Un borde es una cadena que es prefijo y sufijo. Si el problema pide bordes de \textbf{todas} las subcadenas, no
alcanza KMP sobre la cadena completa. Cada borde no vacio de longitud \(b\) corresponde a dos ocurrencias iguales
de longitud \(b\) en los extremos de alguna subcadena. Por tanto:
\[
ans=\frac{n(n+1)}2+\sum_{i<j}\operatorname{lcp}(i,j),
\]
donde el primer termino cuenta el borde vacio si el statement lo incluye.
La suma de LCP de todos los pares de sufijos equivale a la suma de minimos de todos los subarreglos del arreglo
\(\LCP\), calculable con pila monotona en \(\bigO(n)\) despues de construir \(\SA+\LCP\).

\newpage

\section{Arboles estaticos: LCA, HLD, Imperial Roads y Fools and Roads}

\subsection{LCA con binary lifting}
\[
up[v][j]=\text{ancestro de }v\text{ a distancia }2^j.
\]
Si se necesitan maximos en caminos:
\[
mx[v][j]=\max\text{ peso de arista al subir de }v\text{ hacia }up[v][j].
\]
Para responder \(maxEdge(u,v)\), se igualan profundidades acumulando maximos; luego se suben ambos nodos
desde potencias grandes hacia pequenas hasta llegar al LCA.

\[
T_{\text{pre}}=\bigO(n\log n),\qquad T_{\text{query}}=\bigO(\log n),\qquad M=\bigO(n\log n).
\]

\subsection{HLD}
HLD marca como pesado al hijo de mayor subarbol. Toda arista no pesada es liviana. Si se baja por una arista
liviana, el tamano del subarbol se reduce al menos a la mitad; por eso cualquier camino raiz-nodo cruza
\(\bigO(\log n)\) aristas livianas. Entonces un camino \(u\to v\) se descompone en \(\bigO(\log n)\) segmentos
contiguos sobre un arreglo base.

\boxhead{Cuando usar HLD}
Consultas o updates sobre caminos de un arbol fijo: suma, maximo, minimo, pintar, agregar valor, comparar
segmentos, etc. Con un segment tree, cada segmento cuesta \(\bigO(\log n)\), por lo que:
\[
T_{\text{path}}=\bigO(\log^2 n).
\]

\subsection{Imperial Roads}
Problema tipo: grafo conectado ponderado. Cada consulta obliga a incluir una arista \(e=(u,v,w)\) en el arbol
de costo minimo que conecta todos los vertices.

Algoritmo:
\begin{enumerate}
\item Construir MST \(T\) con Kruskal. Sea \(W\) su peso.
\item Preprocesar \(T\) para \(maxEdge(u,v)\) con LCA o HLD.
\item Para query \(e=(u,v,w)\):
\[
ans=W+w-\operatorname{maxEdge}_T(u,v).
\]
\end{enumerate}

Correctitud: al agregar \(e\) al MST aparece un unico ciclo formado por \(e\) y el camino \(u\to v\) en \(T\).
Para conservar un arbol que incluya \(e\), hay que eliminar una arista de ese ciclo. El costo se minimiza eliminando
la arista de mayor peso del camino original. Esta es una aplicacion directa de la propiedad de ciclo del MST.

Complejidad:
\[
T=\bigO(m\log m+n\log n+q\log n),\qquad M=\bigO(n\log n+m).
\]

\subsection{Fools and Roads}
Problema tipo: en un arbol, muchas consultas \((u,v)\). Para cada arista, contar cuantas rutas \(u\to v\) la usan.

Algoritmo de diferencia en arbol:
\[
cnt[u]++,\qquad cnt[v]++,\qquad cnt[\lca(u,v)]-=2.
\]
Luego se hace DFS postorder:
\[
sub[x]=cnt[x]+\sum_{hijo\ y} sub[y].
\]
El valor \(sub[x]\) es la cantidad de caminos que usan la arista \((parent[x],x)\).

Correctitud: cada query agrega dos unidades de flujo en los extremos. Ese flujo sube por el arbol y se cancela en
el LCA. Por tanto, una arista recibe flujo exactamente cuando esta en el camino entre los extremos.

Complejidad:
\[
T=\bigO(n\log n+q\log n+n),\qquad M=\bigO(n\log n).
\]

\newpage

\section{Dynamic Graphs: rollback DSU y estructuras dinamicas}

\subsection{Offline dynamic connectivity}
Si hay operaciones de agregar y borrar aristas, pero conocemos toda la secuencia antes de responder, podemos
resolver offline.

\boxhead{Transformacion clave} Cada arista tiene intervalos de vida \([l,r]\) en los que esta activa. Insertamos
ese intervalo en los nodos de un segment tree sobre el tiempo que cubren exactamente \([l,r]\).

\begin{lstlisting}
dfs(node):
  snap = DSU.snapshot()
  for edge in node.edges:
    DSU.union(edge.u, edge.v)
  if node is leaf:
    answer queries at that time
  else:
    dfs(left); dfs(right)
  DSU.rollback(snap)
\end{lstlisting}

\subsection{Rollback DSU}
Un DSU con rollback guarda en una pila los cambios hechos por cada union. Normalmente no usa path compression,
porque comprimir caminos modifica muchos padres y seria dificil deshacer. Usa union by size/rank.

Si ademas se pide suma de componente, cada raiz guarda:
\[
sum[root]=\sum_{v\in component(root)} value[v].
\]
Al unir componentes, se suma. Para rollback, se guardan los valores antiguos de \(parent\), \(size\) y \(sum\).

\subsection{Correctitud}
Un intervalo de vida \([l,r]\) se descompone en \(\bigO(\log q)\) nodos del segment tree. Para una hoja \(t\),
el camino desde la raiz hasta esa hoja contiene exactamente los nodos cuyos intervalos cubren \(t\). Por eso, al
llegar a la hoja, el DSU contiene exactamente las aristas activas en el tiempo \(t\). El rollback restaura el estado
antes de explorar otra rama.

\subsection{Complejidad}
Si hay \(m\) intervalos de vida y \(q\) tiempos:
\[
T=\bigO(m\log q\cdot T_{\text{union}}),\qquad M=\bigO(m\log q+n).
\]
Con union by size sin path compression, \(T_{\text{union}}\) es \(\bigO(\log n)\) en una cota simple, o casi constante
en practica. Lo importante para examen es justificar el factor \(\log q\) por el segment tree de tiempo.

\subsection{Euler-tour tree vs link-cut tree}
\begin{center}
\begin{tabular}{p{0.28\linewidth}p{0.30\linewidth}p{0.30\linewidth}}
\toprule
Caso & Estructura & Idea \\
\midrule
Bosque dinamico, conectividad, size, suma de componente & Euler-tour tree & tour Euler en BST balanceado; link/cut son split/merge \\
Bosque dinamico, max/sum/min en camino & Link-cut tree & preferred paths en splay; \texttt{access} expone un camino \\
Grafo general offline & Rollback DSU & intervalos de vida sobre tiempo \\
Solo inserciones & DSU normal & componentes monotonicamente se unen \\
Solo borrados & Procesar al reves & los borrados se vuelven inserciones \\
\bottomrule
\end{tabular}
\end{center}

\boxhead{Link-cut tree en una frase} \(\texttt{access}(v)\) reorganiza los preferred paths para que el camino
desde la raiz representada hasta \(v\) quede en un splay, donde el agregado del camino esta en la raiz.

\boxhead{Euler-tour tree en una frase} Representa cada arbol por la secuencia de su tour Euler guardada en un
BST balanceado. Si dos vertices pertenecen al mismo BST, estan conectados.

\newpage

\section{Integer structures}

\subsection{Modelo y problema predecessor}
Se mantiene un conjunto \(S\) de \(n\) enteros de \(w\) bits. El universo es:
\[
U=\{0,1,\ldots,2^w-1\},\qquad u=2^w.
\]
Operaciones:
\[
insert(x),\ delete(x),\ member(x),\ pred(x),\ succ(x).
\]
En modelo de comparaciones, predecessor cuesta \(\bigO(\log n)\) con un BST. En word RAM se usan bits y el
universo para mejorar.

\subsection{van Emde Boas}
vEB divide el universo en coordenadas high/low:
\[
x=(high(x),low(x)),
\]
donde cada parte tiene tamano \(\sqrt u\). Guarda:
\begin{itemize}
\item \(min,max\);
\item clusters para las partes altas;
\item un summary que indica que clusters no estan vacios.
\end{itemize}
La recurrencia es:
\[
T(u)=T(\sqrt u)+\bigO(1).
\]
Como \(u=2^w\), cada llamada divide el numero de bits por \(2\):
\[
T(w)=T(w/2)+\bigO(1)=\bigO(\log w)=\bigO(\log\log u).
\]
El problema es el espacio:
\[
M=\Theta(u)
\]
en la version basica.

\subsection{Y-fast trees}
Y-fast reduce el espacio a lineal. La idea es:
\begin{itemize}
\item mantener representantes de buckets en una estructura tipo X-fast;
\item cada bucket tiene tamano \(\bigO(w)\) y se guarda con una estructura balanceada;
\item split/merge mantiene buckets de tamano controlado.
\end{itemize}
Resultado teorico:
\[
T_{\text{op}}=\bigO(\log w)=\bigO(\log\log u) \text{ esperado/amortizado},\qquad M=\bigO(n).
\]
Si el universo es enorme y \(n\) es pequeno, Y-fast es mejor respuesta que vEB basico por espacio.

\subsection{Fusion trees}
Fusion trees usan branching alto \(B=w^{1/5}\). Dentro de un nodo se guardan llaves y se extraen bits criticos
del trie conceptual para formar sketches. Esos sketches se empaquetan en palabras y se comparan en paralelo.
La altura baja a:
\[
\bigO(\log_B n)=\bigO(\log_w n).
\]
En examen, mencionar que depende de operaciones word RAM fuertes; algunas versiones usan multiplicacion, lo
cual importa si el modelo es AC0 RAM.

\subsection{Ordenamiento de enteros}
La cota \(\Omega(n\log n)\) aplica al modelo de comparaciones. Con enteros se pueden mirar digitos.

\boxhead{Counting sort}
Si las llaves estan en universo \(U_d\):
\[
T=\bigO(n+U_d),\qquad M=\bigO(U_d).
\]

\boxhead{Radix sort}
Partir cada llave en digitos de \(b\) bits. Cada pasada usa counting sort estable:
\[
T=\bigO\left((n+2^b)\left\lceil\frac{w}{b}\right\rceil\right).
\]
Si \(b=\Theta(\log n)\), entonces \(2^b=\Theta(n)\) y:
\[
T=\bigO\left(n\frac{w}{\log n}\right).
\]
Si \(w=\bigO(\log n)\), se obtiene tiempo lineal.

\boxhead{Signature sort}
Para enteros grandes, se reducen llaves a firmas/hashes de \(\bigO(\log n)\) bits por niveles, y luego se aplica
packed/radix sort. La idea teorica es reducir los bits relevantes sin perder el orden.

\newpage

\section{Patrones de respuesta para problemas Racso}

\subsection{Si aparece un problema de strings}
Primero intenta ubicarlo en uno de estos moldes:
\begin{itemize}
\item Buscar o contar patrones: intervalo en suffix array.
\item Numero de subcadenas distintas: \(\frac{n(n+1)}2-\sum \LCP\).
\item Longest common substring: SA de \(A\#B\) y max LCP entre origen distinto.
\item Ordenar/comparar subcadenas: \(lcp(i,j)\) con RMQ sobre LCP.
\item Muchos patrones en un texto: Aho-Corasick para ocurrencias, SA si el texto fijo es lo central.
\item Periodicidad: igualdad de bloques via LCP; si es avanzado, mencionar runs.
\end{itemize}

\subsection{Si aparece un problema de arboles}
\begin{itemize}
\item Una ruta especifica \(u\to v\): LCA o HLD.
\item Muchas rutas y contar aristas usadas: diferencia en arbol.
\item Maximo/minimo en path de un arbol fijo: HLD o binary lifting augmentado.
\item Contar pares de nodos por distancia: centroid decomposition.
\item Subarboles: Euler tour timestamps.
\end{itemize}

\subsection{Si aparece dynamic graph}
\begin{itemize}
\item Solo inserciones: DSU.
\item Solo borrados: procesar al reves.
\item Inserciones y borrados offline: segment tree sobre tiempo + rollback DSU.
\item Bosque dinamico online con path aggregates: link-cut tree.
\item Bosque dinamico online con componente/subtree aggregates: Euler-tour tree.
\item Grafo general online: spanning forests por niveles + Euler-tour trees; es mas teorico.
\end{itemize}

\subsection{Si aparece integer}
Siempre empieza definiendo \(n,w,u=2^w\) y el modelo. Luego:
\begin{itemize}
\item Predecessor con universo moderado: vEB, \(\bigO(\log\log u)\), espacio \(\Theta(u)\).
\item Predecessor con espacio lineal: Y-fast, \(\bigO(\log w)\), espacio \(\bigO(n)\).
\item Predecessor teorico rapido: Fusion tree, \(\bigO(\log_w n)\).
\item Sorting de enteros: counting/radix si \(w\) es pequeno; signature si \(w\) es grande.
\end{itemize}

\section{Errores que quitan puntos}
\begin{itemize}
\item Confundir rango cerrado \([l,r]\) con medio abierto \([l,r)\) en RMQ.
\item Decir ``uso suffix array'' sin explicar intervalo contiguo o LCP.
\item En HLD con valores en aristas, olvidar excluir la posicion del LCA.
\item Usar path compression en rollback DSU.
\item Para Imperial Roads, olvidar que se fuerza una arista y se quita la maxima del ciclo.
\item Para bordes, resolver solo bordes de la cadena completa cuando el problema pide todas las subcadenas.
\item No mencionar memoria. En integer, vEB tiene espacio \(\Theta(u)\), no lineal.
\item Usar int para cantidades que pueden ser \(\Theta(n^2)\): subcadenas, pares de nodos, sumas de LCP.
\end{itemize}

\section{Resumen final de formulas y complejidades}
Esta tabla es para repasar en los ultimos minutos antes de entrar.

\begin{center}
\begin{tabular}{p{0.25\linewidth}p{0.42\linewidth}p{0.22\linewidth}}
\toprule
Tema & Formula / idea que debes recordar & Costo \\
\midrule
Sparse table &
\(st[k][i]=op(st[k-1][i],st[k-1][i+2^{k-1}])\). Query con dos bloques de tamano \(2^k\). &
Build \(\bigO(n\log n)\), query \(\bigO(1)\), memoria \(\bigO(n\log n)\). \\
\addlinespace
LCA binary lifting &
\(up[v][j]\) ancestro \(2^j\). Para max edge, guardar \(mx[v][j]\). &
Pre \(\bigO(n\log n)\), query \(\bigO(\log n)\). \\
\addlinespace
HLD &
Light edge reduce subarbol al menos a la mitad. Path \(u\to v\) son \(\bigO(\log n)\) segmentos. &
Pre \(\bigO(n)\), path con segtree \(\bigO(\log^2 n)\). \\
\addlinespace
Suffix array &
\(\rankv[i]\) es posicion del sufijo \(i\) en \(\SA\). \(\LCP[p]=lcp(\SA[p-1],\SA[p])\). &
SA \(\bigO(n\log n)\) o \(\bigO(n\log^2 n)\), Kasai \(\bigO(n)\). \\
\addlinespace
LCP arbitrario &
Si \(a=\rankv[i]<b=\rankv[j]\), \(lcp(i,j)=\min \LCP[a+1..b]\). &
RMQ build \(\bigO(n\log n)\), query \(\bigO(1)\). \\
\addlinespace
Subcadenas distintas &
\(\#=\frac{n(n+1)}2-\sum_i \LCP[i]\). &
\(\bigO(n)\) despues de SA+LCP. \\
\addlinespace
Imperial Roads &
\(ans=W+w-\operatorname{maxEdge}_T(u,v)\). &
Kruskal \(\bigO(m\log m)\), query \(\bigO(\log n)\). \\
\addlinespace
Fools and Roads &
\(cnt[u]++,cnt[v]++,cnt[\lca(u,v)]-=2\); postorder acumula. &
\(\bigO((n+q)\log n)\). \\
\addlinespace
Centroid paths \(\le C\) &
En centroid: contar pares de distancias con two pointers; restar pares del mismo subarbol. &
\(\bigO(n\log^2 n)\), memoria \(\bigO(n)\). \\
\addlinespace
Rollback DSU &
Arista activa en intervalo \([l,r]\); guardar en segment tree de tiempo; DFS con snapshot/rollback. &
\(\bigO(m\log q)\) uniones, memoria \(\bigO(m\log q+n)\). \\
\addlinespace
vEB &
\(T(u)=T(\sqrt u)+\bigO(1)=\bigO(\log\log u)\). &
Tiempo \(\bigO(\log\log u)\), memoria \(\Theta(u)\). \\
\addlinespace
Radix sort &
\(\bigO((n+2^b)\lceil w/b\rceil)\), elegir \(b=\Theta(\log n)\). &
\(\bigO(nw/\log n)\), lineal si \(w=\bigO(\log n)\). \\
\bottomrule
\end{tabular}
\end{center}

\boxhead{Frase final para casi toda respuesta}
``La estructura es correcta porque mantiene exactamente la informacion que la consulta necesita: minimo de
rangos canonicos, prefijo comun de sufijos, maximo en un camino, componentes activas en un tiempo, o bits
relevantes del universo. El algoritmo no enumera objetos cuadraticos; los comprime en intervalos, caminos,
clusters, buckets o descomposiciones balanceadas.''

\end{document}
